\subsection{ESPRIT (Estimation of Signal Parameters via Rotational Invariance Techniques)}\label{sec:doa:esprit}

El algoritmo \textbf{ESPRIT} es una alternativa a MUSIC que evita el barrido angular.  
En lugar de calcular un espectro y buscar picos, se apoya en la 
\textbf{invariancia rotacional} entre subarreglos de un ULA (Uniform Linear Array).  
Esto permite estimar directamente las direcciones de arribo resolviendo un problema algebraico \cite{Gupta2015}.

\subsubsection*{Idea general}
\begin{itemize}
	\item Se divide el arreglo de $M$ sensores en dos subarreglos idénticos, 
	desplazados por una distancia fija $\Delta x$.  
	\item La señal en el segundo subarreglo es equivalente a la del primero, 
	multiplicada por un factor de fase dependiente del ángulo de arribo.  
	\item Esta relación se modela como una \textbf{invariancia rotacional}, 
	lo que permite estimar los ángulos a partir de los autovalores de una matriz reducida.  
\end{itemize}

Mientras que MUSIC “escanea” ángulo por ángulo para detectar picos,  
ESPRIT utiliza propiedades algebraicas del arreglo y evita la búsqueda explícita.

\subsubsection*{Características}
\begin{itemize}
	\item No requiere barrido angular, reduciendo la complejidad frente a MUSIC.  
	\item Menor tiempo de cálculo y menos sensible a discretización de ángulos.  
	\item Exige que el arreglo posea subestructuras desplazadas (ej. ULA).  
	\item Sensible a errores de calibración o desviaciones geométricas.  
\end{itemize}

\subsubsection*{Formulación matemática}

Dividiendo el arreglo en dos subarreglos, las señales medidas se modelan como:
\begin{equation}
	\mathbf{u}_0(t) = A\, \mathbf{s}(t) + \mathbf{n}_0(t), 
	\qquad 
	\mathbf{u}_1(t) = A \Phi \mathbf{s}(t) + \mathbf{n}_1(t),
\end{equation}
donde $A$ es la matriz de steering del primer subarreglo, 
y $\Phi$ una matriz diagonal de factores de fase:
\[
\Phi = \mathrm{diag}\left(e^{j\gamma_1}, e^{j\gamma_2}, \dots, e^{j\gamma_L}\right), 
\qquad \gamma_l = \frac{2\pi}{\lambda}\, \Delta x \cos\varphi_l.
\]

A partir del subespacio de señal $V_s$, se extraen dos bloques $V_0$ y $V_1$ 
correspondientes a cada subarreglo.  
La relación de invariancia se expresa como:
\begin{equation}
	V_1 \approx V_0 \Psi,
\end{equation}
donde $\Psi$ es una matriz de rotación que contiene la información angular.

Los autovalores de $\Psi$ son:
\[
\exp(j \gamma_l) = e^{j \tfrac{2\pi}{\lambda}\, \Delta x \cos\varphi_l},
\]
a partir de los cuales se obtiene directamente la estimación de los ángulos:
\begin{equation}
	\hat{\varphi}_l = \cos^{-1}\!\left( \frac{\lambda}{2\pi \Delta x} \, \angle(e^{j \gamma_l}) \right), 
	\quad l=0,\dots,L-1.
\end{equation}
